% =========================================================
% Sigma-Algebras
% =========================================================

\subsection{Sigma-Algebras}

\begin{exposition}
Set algebras control finite Boolean manipulation.  Measure theory needs one
more closure rule: countable unions.  A sigma-algebra is the structure that
keeps the finite Boolean workspace and adds closure under countable
aggregation.
\end{exposition}

\begin{definitionbox}{Definition (Sigma-Algebra)}
\begin{definition}[Sigma-Algebra]\label{def:sigma-algebra}
Let \(X\) be a set.  A \emph{sigma-algebra} on \(X\) is a family
\(\Sigma\subseteq\mathcal P(X)\) such that:
\begin{enumerate}[label=(\roman*)]
\item \(\Sigma\) is a set algebra on \(X\);
\item whenever \(A_n\in\Sigma\) for every \(n\in\mathbb N\), the countable
      union \(\bigcup_{n\in\mathbb N}A_n\) belongs to \(\Sigma\).
\end{enumerate}
\end{definition}
\end{definitionbox}

\begin{remark*}[Standard quantified statement]
\[
\Sigma\text{ is a sigma-algebra on }X
\iff
\Sigma\text{ is a set algebra on }X
\ \wedge\
\forall (A_n)_{n\in\mathbb N}\subseteq\Sigma,\ 
\bigcup_{n\in\mathbb N}A_n\in\Sigma.
\]
\end{remark*}

\begin{remark*}[Predicate reading]
\[
  \Sigma\text{ is a sigma-algebra on }X
  \iff
  \Sigma\text{ is a set algebra on }X
  \wedge
  \forall (A_n)_{n\in\mathbb N}\subseteq\Sigma,\ 
  \bigcup_{n\in\mathbb N}A_n\in\Sigma.
\]
\end{remark*}

\begin{remark*}[Interpretation]
A sigma-algebra is the first family structure that supports limiting
operations.  The countable-union axiom is what later makes measurable sets
stable under countable decompositions, approximations, and limiting
constructions.
\end{remark*}

\begin{remark*}[Exposition]
The chapter is deliberately not asking for closure under all unions.  Arbitrary
union closure belongs to topology.  Sigma-algebras require countable union
closure because measure-theoretic additivity and limiting arguments are
countable.
\end{remark*}

\LeanFormalizes{def:sigma-algebra}{lra-lean}{LRA.VolumeI.Set.Families}{SigmaAlgebra}{pending}

\begin{dependencies}
\begin{itemize}
  \item \hyperref[def:set-algebra]{Set Algebra}
  \item \hyperref[def:closed-under-set-operation-vocabulary]{Closed Under Set Operations}
  \item \hyperref[def:indexed-family]{Indexed Family}
  \item \hyperref[def:indexed-union]{Indexed Union}
\end{itemize}
\end{dependencies}

\begin{theorembox}{Theorem (Sigma-Algebra Closure Consequences)}
\begin{theorem}[Sigma-Algebra Closure Consequences]\label{thm:sigma-algebra-closure-consequences}
Let \(\Sigma\) be a sigma-algebra on \(X\).  Then:
\begin{enumerate}[label=(\roman*)]
\item \(\Sigma\) is a set algebra on \(X\);
\item \(\varnothing\in\Sigma\);
\item \(\Sigma\) is closed under complements;
\item \(\Sigma\) is closed under finite unions;
\item \(\Sigma\) is closed under finite intersections;
\item \(\Sigma\) is closed under countable unions;
\item \(\Sigma\) is closed under countable intersections.
\end{enumerate}
\hyperref[prf:sigma-algebra-closure-consequences]{Go to proof.}
\end{theorem}
\end{theorembox}

\begin{remark*}[Standard quantified statement]
\[
\Sigma\text{ is a sigma-algebra on }X
\to
\Sigma\text{ is closed under finite Boolean operations, countable unions, and countable intersections}.
\]
\end{remark*}

\begin{remark*}[Interpretation]
This theorem is the working license for measure-theoretic set manipulation.
Once a family is known to be a sigma-algebra, finite Boolean manipulations and
countable union/intersection manipulations stay inside the family.
\end{remark*}

\begin{remark*}[Exposition]
Closure under countable intersections is not a separate axiom here.  It follows
from complement closure, countable-union closure, and indexed De Morgan.
\end{remark*}

\LeanFormalizes{thm:sigma-algebra-closure-consequences}{lra-lean}{LRA.VolumeI.Set.Families}{SigmaAlgebraIsSetAlgebra}{pending}
\LeanFormalizes{thm:sigma-algebra-closure-consequences}{lra-lean}{LRA.VolumeI.Set.Families}{SigmaAlgebraContainsEmpty}{pending}
\LeanFormalizes{thm:sigma-algebra-closure-consequences}{lra-lean}{LRA.VolumeI.Set.Families}{SigmaAlgebraClosedUnderComplements}{pending}
\LeanFormalizes{thm:sigma-algebra-closure-consequences}{lra-lean}{LRA.VolumeI.Set.Families}{SigmaAlgebraClosedUnderFiniteUnions}{pending}
\LeanFormalizes{thm:sigma-algebra-closure-consequences}{lra-lean}{LRA.VolumeI.Set.Families}{SigmaAlgebraClosedUnderFiniteIntersections}{pending}
\LeanFormalizes{thm:sigma-algebra-closure-consequences}{lra-lean}{LRA.VolumeI.Set.Families}{SigmaAlgebraClosedUnderCountableUnions}{pending}
\LeanFormalizes{thm:sigma-algebra-closure-consequences}{lra-lean}{LRA.VolumeI.Set.Families}{SigmaAlgebraClosedUnderCountableIntersections}{pending}

\begin{dependencies}
\begin{itemize}
  \item \hyperref[def:sigma-algebra]{Sigma-Algebra}
  \item \hyperref[thm:set-algebra-closure-consequences]{Set Algebra Closure Consequences}
  \item \hyperref[thm:indexed-de-morgan]{De Morgan's Laws for Indexed Families}
  \item \hyperref[def:closed-under-set-operation-vocabulary]{Closed Under Set Operations}
  \item \hyperref[def:indexed-union]{Indexed Union}
  \item \hyperref[def:indexed-intersection]{Indexed Intersection}
\end{itemize}
\end{dependencies}

\begin{theorembox}{Theorem (Sigma-Algebras Are Closed Under Intersection)}
\begin{theorem}[Sigma-Algebras Are Closed Under Intersection]\label{thm:sigma-algebras-closed-under-intersection}
Let \(\mathfrak S\) be a family of sigma-algebras on \(X\).  Then
\[
  \bigcap_{\Sigma\in\mathfrak S}\Sigma
\]
is a sigma-algebra on \(X\).
\hyperref[prf:sigma-algebras-closed-under-intersection]{Go to proof.}
\end{theorem}
\end{theorembox}

\begin{remark*}[Standard quantified statement]
\[
\forall\mathfrak S,\quad
\bigl(\forall\Sigma\in\mathfrak S,\ \Sigma\text{ is a sigma-algebra on }X\bigr)
\to
\bigcap_{\Sigma\in\mathfrak S}\Sigma\text{ is a sigma-algebra on }X.
\]
\end{remark*}

\begin{remark*}[Interpretation]
This theorem is the structural reason generated sigma-algebras exist.  To build
the smallest sigma-algebra containing a family \(\mathcal E\), intersect all
sigma-algebras that contain \(\mathcal E\).  The intersection remains a
sigma-algebra.
\end{remark*}

\begin{remark*}[Exposition]
If \(\mathfrak S\) is empty, the intersection is interpreted as the whole
power set \(\mathcal P(X)\), which is a sigma-algebra.  In generated
sigma-algebra applications, the indexing family is nonempty anyway because
\(\mathcal P(X)\) itself contains every starting family.
\end{remark*}

\LeanFormalizes{thm:sigma-algebras-closed-under-intersection}{lra-lean}{LRA.VolumeI.Set.Families}{SigmaAlgebrasClosedUnderIntersection}{pending}

\begin{dependencies}
\begin{itemize}
  \item \hyperref[def:sigma-algebra]{Sigma-Algebra}
  \item \hyperref[thm:sigma-algebra-closure-consequences]{Sigma-Algebra Closure Consequences}
  \item \hyperref[thm:closure-rule-intersection-stability]{Closure Rule Intersection Stability}
  \item \hyperref[def:collection-system]{Collection System}
  \item \hyperref[def:indexed-intersection]{Indexed Intersection}
\end{itemize}
\end{dependencies}

\begin{definitionbox}{Definition (Generated Sigma-Algebra)}
\begin{definition}[Generated Sigma-Algebra]\label{def:generated-sigma-algebra}
Let \(\mathcal E\subseteq\mathcal P(X)\).  The \emph{sigma-algebra generated
by \(\mathcal E\)} is
\[
  \sigma(\mathcal E)
  =
  \bigcap\{\Sigma\subseteq\mathcal P(X):
    \Sigma\text{ is a sigma-algebra on }X
    \text{ and }\mathcal E\subseteq\Sigma\}.
\]
\end{definition}
\end{definitionbox}

\begin{remark*}[Standard quantified statement]
\[
\sigma(\mathcal E)
=
\bigcap\{\Sigma\subseteq\mathcal P(X):
\Sigma\text{ is a sigma-algebra on }X
\text{ and }\mathcal E\subseteq\Sigma\}.
\]
\end{remark*}

\begin{remark*}[Interpretation]
\(\sigma(\mathcal E)\) is the smallest sigma-algebra that contains the
starting family \(\mathcal E\).  It is the measure-theoretic analogue of
``close this data under the allowed operations, but add nothing unnecessary.''
\end{remark*}

\begin{remark*}[Exposition]
The family of candidate sigma-algebras is nonempty because
\(\mathcal P(X)\) is a sigma-algebra containing every
\(\mathcal E\subseteq\mathcal P(X)\).  Therefore the intersection construction
has at least one candidate to intersect.
\end{remark*}

\LeanFormalizes{def:generated-sigma-algebra}{lra-lean}{LRA.VolumeI.Set.Families}{SigmaAlgebraSystem}{pending}
\LeanFormalizes{def:generated-sigma-algebra}{lra-lean}{LRA.VolumeI.Set.Families}{GeneratedSigmaAlgebra}{pending}

\begin{dependencies}
\begin{itemize}
  \item \hyperref[def:sigma-algebra]{Sigma-Algebra}
  \item \hyperref[def:generated-collection]{Generated Collection}
  \item \hyperref[thm:sigma-algebras-closed-under-intersection]{Sigma-Algebras Are Closed Under Intersection}
  \item \hyperref[def:subset]{Subset}
\end{itemize}
\end{dependencies}

\begin{theorembox}{Theorem (Generated Sigma-Algebra Closure Laws)}
\begin{theorem}[Generated Sigma-Algebra Closure Laws]\label{thm:generated-sigma-algebra-closure-laws}
Let \(\mathcal E,\mathcal F\subseteq\mathcal P(X)\).  Then:
\begin{enumerate}[label=(\roman*)]
\item \(\mathcal E\subseteq\sigma(\mathcal E)\);
\item if \(\mathcal E\subseteq\mathcal F\), then
      \(\sigma(\mathcal E)\subseteq\sigma(\mathcal F)\);
\item \(\sigma(\mathcal E)\) is a sigma-algebra on \(X\);
\item \(\sigma(\sigma(\mathcal E))=\sigma(\mathcal E)\).
\end{enumerate}
\hyperref[prf:generated-sigma-algebra-closure-laws]{Go to proof.}
\end{theorem}
\end{theorembox}

\begin{remark*}[Standard quantified statement]
\[
\mathcal E\subseteq\sigma(\mathcal E),\quad
\mathcal E\subseteq\mathcal F\to\sigma(\mathcal E)\subseteq\sigma(\mathcal F),\quad
\sigma(\mathcal E)\text{ is a sigma-algebra on }X,\quad
\sigma(\sigma(\mathcal E))=\sigma(\mathcal E).
\]
\end{remark*}

\begin{remark*}[Interpretation]
Generated sigma-algebra formation is a closure operator on families of subsets:
it contains the starting data, respects inclusion of starting data, lands in
sigma-algebras, and stabilizes after one application.
\end{remark*}

\LeanFormalizes{thm:generated-sigma-algebra-closure-laws}{lra-lean}{LRA.VolumeI.Set.Families}{GeneratedSigmaAlgebraCandidateNonempty}{pending}
\LeanFormalizes{thm:generated-sigma-algebra-closure-laws}{lra-lean}{LRA.VolumeI.Set.Families}{GeneratedSigmaAlgebraExtensive}{pending}
\LeanFormalizes{thm:generated-sigma-algebra-closure-laws}{lra-lean}{LRA.VolumeI.Set.Families}{GeneratedSigmaAlgebraMonotone}{pending}
\LeanFormalizes{thm:generated-sigma-algebra-closure-laws}{lra-lean}{LRA.VolumeI.Set.Families}{GeneratedSigmaAlgebraIsSigmaAlgebra}{pending}
\LeanFormalizes{thm:generated-sigma-algebra-closure-laws}{lra-lean}{LRA.VolumeI.Set.Families}{GeneratedSigmaAlgebraIdempotent}{pending}

\begin{dependencies}
\begin{itemize}
  \item \hyperref[def:generated-sigma-algebra]{Generated Sigma-Algebra}
  \item \hyperref[thm:generated-collection-closure-laws]{Generated Collection Closure Laws}
  \item \hyperref[thm:sigma-algebras-closed-under-intersection]{Sigma-Algebras Are Closed Under Intersection}
  \item \hyperref[def:subset]{Subset}
\end{itemize}
\end{dependencies}
